% English translation of checked_pages/esq_119_051.tex; source: cote 119, site 51, PDF page 52, typescript page 49.
% AI-translated by GPT-6 Astra (OpenAI, ChatGPT Pro mode).
\scanpage{119/51}
\mspage{49}

quite often, considerable psychological obstacles. That is why,
in certain respects and increasingly in recent years,
my teaching activities have often acted as a burden,
rather than as a stimulus to the development of geometric
reflection that is at all advanced, or merely delicate.

9. This seems an opportune occasion to make a brief
appraisal of my teaching activity since 1970, that is, since
it has taken place in a university setting. This contact
with a very different reality has taught me a great
deal, of an order quite different, moreover, from the
merely pedagogical or scientific. This is not the place to
dwell on that subject. I also spoke at the beginning of this report of the
role played by this change of professional milieu in the
renewal of my approach to mathematics, and of my
mathematical centers of interest. On the other hand, when I take stock
of my teaching activity at the level of research properly
speaking, I arrive at a finding of clear and unambiguous failure. In the more
than ten years that this activity has continued year after year within
the same university institution, I have not, at any time, managed
to bring about a place where ``something is happening'' --- where something
``passes'',\tnote{TN-051-01}{The French repeats the verb ``passer'': ``il se passe quelque chose'' means that something is happening, and ``quelque chose passe'' says that something passes among the people. The second literal phrasing retains that verbal echo.} among a group, however small, of people
linked by a common adventure. On two occasions, it is true, around
the years 74 to 76, I had the pleasure and privilege of inspiring
a student to undertake work of scope, pursued with enthusiasm:
Yves Ladegaillerie's work mentioned earlier (par. 3) on
questions of isotopy in dimension 2, and Carlos Contou-Carrère's
(whose mathematical passion had not waited for his encounter
with me to blossom) unpublished work on local and global
Jacobians over general base schemes (part of which
was announced in a note in CR).\tnote{TN-051-02}{CR abbreviates \emph{Comptes Rendus} of the Académie des sciences; see the \href{https://comptes-rendus.academie-sciences.fr/page/history-comptes-rendus/}{journal's historical account}.} Apart from these two cases,
my role over these ten years has been confined to passing on,
as best I could, the rudiments of the mathematician's craft to some
twenty students\unsure{.}\tnote{TN-051-03}{The doubtful dot after ``élèves'' is retained after ``students''.} at the research level, or at least to

\margin{51}
